\section{Расширения полей}

Пусть $F$~--- поле. \emph{Расширением} поля $F$ называется поле, содержащее $F$ в качестве подполя. Иначе говоря,
расширение~--- это $F$-алгебра, кольцо которой является полем. Расширение $E$ поля $F$ является, в частности, векторным
пространством над $F$; его размерность называется \emph{степенью} $E$ над $F$ и обозначается $[E:F]$. Расширение
называется \emph{конечным} (соответственно \emph{квадратичным}, \emph{кубическим} и т.\,д.), если его степень конечна
(соответственно равна $2$, $3$ и т.\,д.). Запись $E/F$ означает <<$E$~--- расширение $F$>>.

Если $E$ и $E'$~--- расширения $F$, то \emph{$F$-гомоморфизмом} $E\to E'$ называется гомоморфизм $\varphi\colon E\to E'$,
для которого $\varphi(c)=c$ при всех $c\in F$; \emph{$F$-изоморфизм}~--- биективный $F$-гомоморфизм. Гомоморфизм полей
всегда инъективен (см.\ после \rf{1.3}), поэтому $F$-гомоморфизм $E\to E'$~--- это в
точности инъективное $F$-линейное отображение, сохраняющее умножение. В частности, если $[E:F]=[E':F]<\infty$, то каждый
$F$-гомоморфизм $E\to E'$ является $F$-изоморфизмом (инъективное линейное отображение пространств одной конечной
размерности сюръективно).

\begin{example}[1.19]
\begin{enumerate}[label=(\alph*),leftmargin=*]
\item Поле $\CC$ имеет степень $2$ над $\R$ (базис $\{1,i\}$).
\item Поле $\R$ имеет бесконечную степень над $\Q$. Действительно, $\Q$ счётно, поэтому всякое конечномерное
$\Q$-пространство, будучи изоморфным $\Q^n$, тоже счётно, а $\R$ несчётно (Кантор: $|\R|=|\mathcal P(\mathbb N)|>|\mathbb N|$, см.\ \rf{6.4d}).
\item Поле гауссовых чисел $\Q(i)=\{a+bi\in\CC\mid a,b\in\Q\}$ имеет степень $2$ над $\Q$ (базис $\{1,i\}$).
\item Поле рациональных функций $F(X)$ имеет бесконечную степень над $F$: уже его подпространство $F[X]$ бесконечномерно
(базис $1,X,X^2,\dots$).
\end{enumerate}
\end{example}

\begin{proposition}[1.20: мультипликативность степеней]
Пусть $L\supset E\supset F$~--- поля. Расширение $L/F$ конечно тогда и только тогда, когда конечны $L/E$ и $E/F$, и в
этом случае
\[
[L:F]=[L:E]\,[E:F].
\]
Точнее, если $(e_i)_{1\le i\le m}$~--- базис $E$ над $F$, а $(l_j)_{1\le j\le n}$~--- базис $L$ над $E$, то
$(e_il_j)_{i,j}$~--- базис $L$ над $F$.
\end{proposition}

\par\noindent\begin{minipage}{\linewidth}
\[
\begin{tikzpicture}[baseline]
\node (L) at (0,3) {$L$};
\node (E) at (0,1.5) {$E$};
\node (F) at (0,0) {$F$};
\draw (L)--(E)--(F);
\draw[decorate,decoration={brace,amplitude=4pt},cold] (0.35,2.85)--(0.35,1.65) node[midway,right=4pt]{$n=[L:E]$, базис $(l_j)$};
\draw[decorate,decoration={brace,amplitude=4pt},cold] (0.35,1.35)--(0.35,0.15) node[midway,right=4pt]{$m=[E:F]$, базис $(e_i)$};
\draw[decorate,decoration={brace,amplitude=5pt,mirror},cnew] (-0.35,2.85)--(-0.35,0.15) node[midway,left=5pt]{$mn$, базис $(e_il_j)$};
\end{tikzpicture}
\]
\figcap{Башня $F\subset E\subset L$: базисы этажей (синие) перемножаются и дают базис всей башни (красный).}
\end{minipage}\par\medskip

\begin{proof}
Если $L$ конечно над $F$, то оно конечно и над $E$ (система, порождающая $L$ над $F$, порождает его и над большим полем
$E$), а $E$, будучи подпространством конечномерного $F$-пространства $L$, конечномерно над $F$.

Пусть теперь $L/E$ и $E/F$ конечны, а $(e_i)$, $(l_j)$~--- базисы, как в формулировке. Покажем, что $(e_il_j)$~--- базис
$L$ над $F$.

\emph{Порождение.} Пусть $\gamma\in L$. Так как $(l_j)$ порождает $L$ над $E$, имеем $\gamma=\sum_j\alpha_jl_j$ с
$\alpha_j\in E$; так как $(e_i)$ порождает $E$ над $F$, имеем $\alpha_j=\sum_ia_{ij}e_i$ с $a_{ij}\in F$. Подставляя,
получаем $\gamma=\sum_{i,j}a_{ij}\,e_il_j$.

\emph{Линейная независимость.} Пусть $\sum_{i,j}a_{ij}e_il_j=0$, $a_{ij}\in F$. Перепишем это как
$\sum_j\bigl(\sum_ia_{ij}e_i\bigr)l_j=0$. Коэффициенты $\sum_ia_{ij}e_i$ лежат в $E$, и линейная независимость $(l_j)$ над
$E$ даёт $\sum_ia_{ij}e_i=0$ при каждом $j$. Линейная независимость $(e_i)$ над $F$ даёт $a_{ij}=0$ при всех $i,j$.
\end{proof}

В частности, если $L/F$ конечно, то $[E:F]$ и $[L:E]$ делят $[L:F]$. Например, если $[L:F]$~--- простое число, то
промежуточных полей, отличных от $F$ и $L$, нет.

\subsection*{Подкольцо, порождённое подмножеством}

Пересечение любого семейства подколец кольца снова является подкольцом (все условия~--- замкнутость относительно сложения,
вычитания, умножения и наличие $1$~--- проверяются поэлементно). Пусть $F$~--- подполе поля $E$ и $S\subset E$.
Пересечение всех подколец $E$, содержащих $F$ и $S$,~--- наименьшее подкольцо $E$, содержащее $F$ и $S$. Оно называется
\emph{подкольцом $E$, порождённым $F$ и $S$} (или \emph{$F$-подалгеброй, порождённой $S$}) и обозначается $F[S]$. Если
$S=\{\alpha_1,\dots,\alpha_n\}$, пишут $F[\alpha_1,\dots,\alpha_n]$. Например, $\CC=\R[\sqrt{-1}]$.

\begin{lemma}[1.21]
Кольцо $F[S]$ состоит из тех элементов $E$, которые записываются конечными суммами вида
\begin{equation}\label{eq:B-poly}
\sum a_{i_1\cdots i_n}\,\alpha_1^{i_1}\cdots\alpha_n^{i_n},\qquad a_{i_1\cdots i_n}\in F,\quad \alpha_k\in S,\quad i_k\in\mathbb N,
\end{equation}
т.\,е.\ значениями многочленов с коэффициентами из $F$ от конечного числа элементов $S$.
\end{lemma}

\begin{proof}
Пусть $R$~--- множество всех таких элементов. Сумма, разность и произведение двух выражений вида \eqref{eq:B-poly}
снова имеют этот вид (объединим наборы участвующих $\alpha_k$ и перемножим многочлены), $1\in R$, и $R$ содержит $F$ и
$S$. Значит, $R$~--- подкольцо, содержащее $F$ и $S$, и $F[S]\subset R$. Обратно, любое подкольцо, содержащее $F$ и $S$,
замкнуто относительно сложения и умножения и потому содержит все суммы \eqref{eq:B-poly}; значит, $R\subset F[S]$.
\end{proof}

\begin{example}[1.22]
Кольцо $\Q[\pi]$, $\pi=3{,}14159\ldots$, состоит из вещественных чисел вида $a_0+a_1\pi+\dots+a_n\pi^n$, $a_i\in\Q$.
Кольцо $\Q[i]$ состоит из комплексных чисел вида $a+bi$, $a,b\in\Q$: все степени $i$ равны $\pm1$ или $\pm i$.
\end{example}

Запись элемента в виде \eqref{eq:B-poly}, вообще говоря, не единственна. Так обстоит дело уже в $\R[i]$: например,
$1+i^2=0=0\cdot i$.

\begin{lemma}[1.23]
Пусть $R$~--- конечная $F$-алгебра. Если $R$~--- область целостности, то $R$~--- поле.
\end{lemma}

\begin{proof}
Пусть $\alpha\in R$, $\alpha\ne0$; нужно найти обратный к $\alpha$ элемент. Отображение $x\mapsto\alpha x\colon R\to R$
$F$-линейно и инъективно (если $\alpha x=0$, то $x=0$, так как в $R$ нет делителей нуля). Инъективное линейное
отображение конечномерного пространства в себя сюръективно, поэтому найдётся $\beta\in R$ с $\alpha\beta=1$.
\end{proof}

В частности, всякое подкольцо конечного расширения $E$ поля $F$, содержащее $F$, является полем: оно конечномерно над $F$
как подпространство $E$ и не имеет делителей нуля как подкольцо поля.

\subsection*{Подполе, порождённое подмножеством}

Пересечение любого семейства подполей поля снова является подполем. Пусть $F$~--- подполе поля $E$ и $S\subset E$.
Пересечение всех подполей $E$, содержащих $F$ и $S$,~--- наименьшее подполе $E$, содержащее $F$ и $S$. Оно называется
\emph{подполем $E$, порождённым $F$ и $S$} (или \emph{порождённым $S$ над $F$}) и обозначается $F(S)$; для
$S=\{\alpha_1,\dots,\alpha_n\}$ пишут $F(\alpha_1,\dots,\alpha_n)$. Поле $F(S)$ совпадает с полем частных $F[S]$ в $E$,
т.\,е.\ с множеством $\{a/b\mid a,b\in F[S],\ b\ne0\}$: это множество~--- подполе $E$, содержащее $F$ и $S$, и оно содержится
в каждом таком подполе (подполе, содержащее $F$ и $S$, содержит $F[S]$ и вместе с ним все частные). Итак,
$F[\alpha_1,\dots,\alpha_n]$ состоит из элементов $E$, выражающихся многочленами от $\alpha_i$ с коэффициентами из $F$, а
$F(\alpha_1,\dots,\alpha_n)$~--- из частных двух таких многочленов (со знаменателем, не равным $0$ в $E$).

По лемме \rf{1.23}, если $F[S]$ конечномерно над $F$, то $F[S]$ уже поле, и $F(S)=F[S]$.

\begin{example}[1.24]
\begin{enumerate}[label=(\alph*),leftmargin=*]
\item Поле $\Q(\pi)$ состоит из вещественных чисел вида $g(\pi)/h(\pi)$, где $g,h\in\Q[X]$, $h\ne0$ (знаменатель $h(\pi)$ не
равен нулю, так как $\pi$ трансцендентно~--- теорема Линдемана, см.\ раздел о трансцендентных числах).
\item Кольцо $\Q[i]$ уже является полем: $(a+bi)^{-1}=\dfrac{a}{a^2+b^2}-\dfrac{b}{a^2+b^2}\,i$ при $(a,b)\ne(0,0)$.
\end{enumerate}
\end{example}

Расширение $E/F$ называется \emph{простым}, если $E=F(\alpha)$ для некоторого $\alpha\in E$. Например, $\Q(\pi)$ и
$\Q[i]$~--- простые расширения $\Q$.

Пусть $F$ и $F'$~--- подполя поля $E$. Пересечение всех подполей $E$, содержащих $F$ и $F'$,~--- наименьшее подполе $E$,
содержащее оба поля. Оно называется \emph{композитом} $F$ и $F'$ в $E$ и обозначается $F\cdot F'$. Это и подполе,
порождённое над $F$ множеством $F'$, и подполе, порождённое над $F'$ множеством $F$:
\[
F(F')=F\cdot F'=F'(F).
\]
Композит определён лишь для двух подполей одного и того же поля. Для двух абстрактных полей он зависит от того, как они
вложены в общее поле: например, в $\CC$ композит $\Q(\sqrt[3]2)\cdot\Q(\omega\sqrt[3]2)$, где $\omega=e^{2\pi i/3}$, имеет
степень $6$ над $\Q$, а $\Q(\sqrt[3]2)\cdot\Q(\sqrt[3]2)=\Q(\sqrt[3]2)$~--- степень $3$, хотя оба сомножителя изоморфны
$\Q[X]/(X^3-2)$. (Первый композит содержит $\omega=\omega\sqrt[3]2/\sqrt[3]2$ и потому равен $\Q(\sqrt[3]2,\omega)$; здесь
$[\Q(\sqrt[3]2):\Q]=3$, так как $X^3-2$ неприводим, \rf{1.16a}(a), а $\omega\notin\Q(\sqrt[3]2)\subset\R$~--- корень
$X^2+X+1$, так что $[\Q(\sqrt[3]2,\omega):\Q(\sqrt[3]2)]=2$; по \rf{1.20} степень равна $3\cdot2$. Степени простых
расширений вычисляются в \rf{1.27c}.)

\subsection*{Построение расширений: факторкольцо $F[X]/(f)$}

Пусть $f\in F[X]$~--- унитарный многочлен степени $m\ge1$, $(f)$~--- порождённый им идеал. Рассмотрим факторкольцо
$F[X]/(f)$ и обозначим через $x$ образ $X$ в нём, т.\,е.\ смежный класс $x=X+(f)$. Композиция $F\to F[X]\to F[X]/(f)$
инъективна (ненулевая константа не делится на $f$, так как $\deg f\ge1$), поэтому $F$ отождествляется с подполем
$F[X]/(f)$.
\begin{enumerate}[label=(\alph*),leftmargin=*]
\item Отображение $P(X)\mapsto P(x)\colon F[X]\to F[x]$~--- это просто проекция на факторкольцо; это $F$-гомоморфизм, и
он переводит $f(X)$ в $0$. Следовательно, $f(x)=0$: \emph{в $F[X]/(f)$ класс $x$~--- корень многочлена $f$}. В
частности, $F[X]/(f)=F[x]$.
\item Каждый элемент $F[X]/(f)$ представляется единственным многочленом $r$ степени $<m$: если $g\in F[X]$ и $g=qf+r$~---
деление с остатком (\rf{1.6}), то $g+(f)=r+(f)$; а если $r_1+(f)=r_2+(f)$ при $\deg r_1,\deg r_2<m$, то $r_1-r_2=qf$, и по
единственности остатка (\rf{1.6}) $r_1=r_2$. Значит, каждый элемент $F[x]$
единственным образом записывается в виде
\begin{equation}\label{eq:B-basis}
a_0+a_1x+\dots+a_{m-1}x^{m-1},\qquad a_i\in F,
\end{equation}
т.\,е.\ $1,x,\dots,x^{m-1}$~--- базис $F[x]$ над $F$, и $\dim_FF[x]=m$.
\item Чтобы сложить два элемента вида \eqref{eq:B-basis}, складывают соответствующие коэффициенты.
\item Чтобы перемножить два элемента вида \eqref{eq:B-basis}, перемножают их как многочлены, а затем с помощью
соотношения $f(x)=0$ выражают одночлены степени $\ge m$ через одночлены меньшей степени: если
$f=X^m+c_{m-1}X^{m-1}+\dots+c_0$, то
\[
x^{m+k}=x^k\cdot x^m=-(c_{m-1}x^{m-1}+\dots+c_0)\,x^k\qquad(k\ge0),
\]
и каждое применение этого правила понижает наибольшую степень; через конечное число шагов получается выражение вида
\eqref{eq:B-basis}. (Это и есть остаток от деления произведения на $f$.)
\item Пусть теперь $f$ неприводим. Тогда каждый ненулевой $\alpha\in F[x]$ обратим, и обратный элемент находится так.
По (b) запишем $\alpha=g(x)$, где $g\in F[X]$, $\deg g\le m-1$, $g\ne0$. Алгоритм Евклида в $F[X]$ (\rf{1.8}) даёт многочлены
$a(X),b(X)$ с
\[
a(X)f(X)+b(X)g(X)=d(X),\qquad d=\gcd(f,g).
\]
Здесь $d=1$: $d$ делит $f$, а делители неприводимого $f$~--- это константы и многочлены вида $cf$, $c\in F^\times$
(\rf{1.6a}); но $d$
делит и $g\ne0$, степень которого меньше $\deg f$, поэтому $d$~--- константа, а $\gcd$ по определению унитарен. Подставив
$x$ вместо $X$ и учитывая $f(x)=0$, получаем
\[
b(x)\,g(x)=1,
\]
т.\,е.\ $b(x)$~--- обратный к $g(x)$ элемент.
\end{enumerate}

Неприводимость в (e) существенна. Если $f=gh$ с $1\le\deg g,\deg h<m$, то $g(x)\ne0$ и $h(x)\ne0$ по (b), но
$g(x)h(x)=f(x)=0$: в $F[X]/(f)$ есть делители нуля, и это не поле. Унитарность $f$ же~--- лишь удобная нормировка: для
$c\in F^\times$ идеалы $(f)$ и $(cf)$ совпадают. Следовательно, для многочлена $f\in F[X]$ степени $m\ge1$ равносильны: $f$
неприводим; $F[X]/(f)$~--- поле; $F[X]/(f)$~--- область целостности (последние два условия равносильны и по \rf{1.23},
так как $\dim_FF[X]/(f)=m<\infty$). Поэтому вопрос о неприводимости $f$ над $F$ (скажем, над расширением исходного поля,
ср.\ упражнение 1-5)~--- это вопрос о том, есть ли делители нуля в $F[X]/(f)$.

Итак, доказано следующее утверждение.

\begin{proposition}[1.25]
Для унитарного неприводимого многочлена $f\in F[X]$ степени $m$ кольцо
\[
F[x]\overset{\mathrm{def}}{=}F[X]/(f)
\]
является полем степени $m$ над $F$, причём $x$~--- корень $f$ в $F[x]$. Вычисления в $F[x]$ сводятся к вычислениям в $F$.
\end{proposition}

Поскольку $F[x]$~--- поле, $F(x)=F[x]$. Обозначение $F[x]$ подчёркивает, что элементы этого поля~--- многочлены от $x$.

Отметим следствие, которое понадобится постоянно: \emph{у каждого непостоянного многочлена $g\in F[X]$ есть корень в
некотором конечном расширении $F$}. Действительно, возьмём унитарный неприводимый делитель $f$ многочлена $g$ (он
существует, так как $F[X]$~--- факториальное кольцо); тогда $x$~--- корень $f$, а значит, и $g$, в поле $F[X]/(f)$ степени
$\deg f\le\deg g$.

\begin{example}[1.26]
Пусть $f=X^2+1\in\R[X]$. Тогда в $\R[x]$:
\begin{itemize}[leftmargin=*]
\item элементы: $a+bx$, $a,b\in\R$;
\item сложение: $(a+bx)+(a'+b'x)=(a+a')+(b+b')x$;
\item умножение: $(a+bx)(a'+b'x)=(aa'-bb')+(ab'+a'b)x$ (так как $x^2=-1$);
\item обратные элементы выписываются сразу: $(a+bx)^{-1}=\dfrac{a-bx}{a^2+b^2}$ при $(a,b)\ne(0,0)$.
\end{itemize}
Обычно пишут $i$ вместо $x$ и $\CC$ вместо $\R[x]$.
\end{example}

\begin{example}[1.27]
Пусть $f=X^3-3X-1\in\Q[X]$. Он неприводим над $\Q$ (\rf{1.12}: рациональных корней у него нет), поэтому $\Q[x]$~--- поле с
базисом $\{1,x,x^2\}$ над $\Q$. Пусть
\[
\beta=x^4+2x^3+3\in\Q[x].
\]
Так как $x^3=3x+1$, имеем $x^4=3x^2+x$, и
\[
\beta=3x^2+x+2(3x+1)+3=3x^2+7x+5.
\]
Найдём $\beta^{-1}$ по способу (e). Алгоритм Евклида для $f$ и $g=3X^2+7X+5$:
\[
X^3-3X-1=\Bigl(\tfrac13X-\tfrac79\Bigr)g+\tfrac{7X+26}{9},\qquad
g=\Bigl(\tfrac{27}{7}X-\tfrac{261}{49}\Bigr)\cdot\tfrac{7X+26}{9}+\tfrac{999}{49}.
\]
Последний остаток~--- ненулевая константа, так что $\gcd(f,g)=1$, как и должно быть. Выражая остатки обратным ходом,
получаем
\[
\bigl(X^3-3X-1\bigr)\Bigl(-\tfrac{7}{37}X+\tfrac{29}{111}\Bigr)+\bigl(3X^2+7X+5\bigr)\Bigl(\tfrac{7}{111}X^2-\tfrac{26}{111}X+\tfrac{28}{111}\Bigr)=1.
\]
Подставляя $x$, находим
\[
\bigl(3x^2+7x+5\bigr)\Bigl(\tfrac{7}{111}x^2-\tfrac{26}{111}x+\tfrac{28}{111}\Bigr)=1,
\]
т.\,е.\ $\beta^{-1}=\frac{7}{111}x^2-\frac{26}{111}x+\frac{28}{111}$. (То же даёт компьютер: в PARI
\texttt{b=Mod(X\^{}4+2*X\^{}3+3,X\^{}3-3*X-1)} и затем \texttt{b\^{}(-1)}; в sympy~--- \texttt{rem} и \texttt{gcdex}.)
\end{example}

\subsection*{Поля корня}

\begin{definition}[1.27a]
Пусть $f\in F[X]$~--- унитарный неприводимый многочлен. Пара $(E,\alpha)$, где $E$~--- расширение $F$ и $\alpha\in E$,
называется \emph{полем корня} многочлена $f$, если $E=F[\alpha]$ и $f(\alpha)=0$.
\end{definition}

Например, пара $(F[X]/(f),x)$~--- поле корня $f$ (\rf{1.25}); назовём её \emph{стандартным} полем корня. Обратите
внимание: в определении требуется $E=F[\alpha]$ (кольцо, порождённое $F$ и $\alpha$), а не $E=F(\alpha)$; то, что кольцо
$F[\alpha]$, порождённое корнем $\alpha$ неприводимого $f$ в \emph{любом} расширении $F$, само оказывается полем, будет
доказано в \rf{1.27c} и объясняется тем, что $\alpha$~--- корень ненулевого многочлена над $F$ (ср.\ \rf{1.22} и случай~1
ниже: для трансцендентного $\pi$ кольцо $\Q[\pi]$ полем не является~--- $\pi$ в нём необратимо).

Здесь важно и то, что $F$~--- \emph{поле}. Корень ненулевого многочлена над полем является корнем и унитарного
многочлена (разделим на старший коэффициент), т.\,е.\ $\alpha$ \emph{цел} над $F$. Целость даёт конечномерность: $F[\alpha]$
порождается как $F$-пространство элементами $1,\alpha,\dots,\alpha^{\deg f-1}$ (старшие степени понижаются с помощью
$f(\alpha)=0$), а конечномерная $F$-алгебра без делителей нуля~--- поле (\rf{1.23}). Над кольцом, не являющимся полем,
одной целости мало: $\sqrt2$~--- корень унитарного многочлена $X^2-2\in\Z[X]$, но $\Z[\sqrt2]$~--- не поле
($\frac12\notin\Z[\sqrt2]$, так как $\sqrt2\notin\Q$).

\begin{proposition}[1.27b]
Пусть $(E,\alpha)$~--- поле корня унитарного неприводимого многочлена $f\in F[X]$ степени $m$. Тогда гомоморфизм
$F$-алгебр $\varphi\colon F[X]\to E$, $g(X)\mapsto g(\alpha)$, сюръективен, его ядро равно $(f)$, и он индуцирует
$F$-изоморфизм
\[
F[x]=F[X]/(f)\xrightarrow{\ \sim\ }E,\qquad x\mapsto\alpha.
\]
Следовательно, каждый элемент $E$ единственным образом записывается в виде $a_0+a_1\alpha+\dots+a_{m-1}\alpha^{m-1}$,
$a_i\in F$, и вычисления в $F[\alpha]$ ведутся по тем же правилам, что в $F[x]$. Если $(E',\alpha')$~--- другое поле корня $f$,
то существует единственный $F$-изоморфизм $E\to E'$, переводящий $\alpha$ в $\alpha'$.
\end{proposition}

\par\noindent\begin{minipage}{\linewidth}
\[
\begin{tikzcd}[column sep=large,row sep=large]
F[X]\arrow[r,"X\,\mapsto\,\alpha",cold]\arrow[d,two heads] & E=F[\alpha]\arrow[r,cnew,"\sim","\alpha\,\mapsto\,\alpha'"'] & E'=F[\alpha']\\
F[X]/(f)\arrow[ur,cnew,"\sim"',"x\,\mapsto\,\alpha"] & &
\end{tikzcd}
\]
\figcap{Поле корня $(E,\alpha)$ изоморфно стандартному $(F[X]/(f),x)$; два поля корня одного многочлена изоморфны
(красные стрелки строятся в доказательстве).}
\end{minipage}\par\medskip

\begin{proof}
\emph{Сюръективность.} Образ $\varphi$~--- это множество значений многочленов от $\alpha$ с коэффициентами из $F$, т.\,е.\
$F[\alpha]$ (\rf{1.21}), а $F[\alpha]=E$ по определению поля корня.

\emph{Ядро.} Идеал $\ker\varphi$ ненулевой: $f\in\ker\varphi$, поскольку $f(\alpha)=0$. По \rf{1.9} $\ker\varphi=(g)$
для единственного унитарного $g$. Значит, $g\mid f$. Делители
неприводимого $f$~--- константы и многочлены $cf$; унитарный $g$ равен либо $1$, либо $f$. Но $g\ne1$, так как
$\varphi(1)=1\ne0$. Итак, $\ker\varphi=(f)$. (Можно рассуждать и так: $\ker\varphi$~--- простой идеал, потому что
$F[X]/\ker\varphi\cong E$~--- область целостности; он собственный и содержит $(f)$, а идеал $(f)$ максимален, так как
$F[X]/(f)$~--- поле (\rf{1.25}); значит, $\ker\varphi=(f)$.)

По теореме о гомоморфизме $\varphi$ индуцирует изоморфизм $F[X]/(f)\to E$, переводящий $x=X+(f)$ в $\varphi(X)=\alpha$
и тождественный на $F$. Единственность записи элементов $E$ переносится с $F[x]$ (п.~(b) построения перед \rf{1.25})
через этот изоморфизм.

Если $(E',\alpha')$~--- другое поле корня, то композиция изоморфизмов $E\to F[x]\to E'$, $\alpha\mapsto x\mapsto\alpha'$,~---
искомый $F$-изоморфизм. Он единствен: $F$-гомоморфизм, переводящий $\alpha$ в $\alpha'$, обязан переводить
$\sum a_i\alpha^i$ в $\sum a_i\alpha'^i$, а такими суммами исчерпывается $E$.
\end{proof}

Иначе говоря: \emph{если поле $E$ порождено над $F$ как кольцо одним элементом $\alpha$, который является корнем
унитарного неприводимого $f\in F[X]$, то $E\cong F[X]/(f)$ как $F$-алгебры}. Поле корня часто обозначают просто $F[\alpha]$.

\subsection*{Алгебраические и трансцендентные элементы}

Пусть $F$~--- поле, а $E$~--- область целостности, содержащая $F$ в качестве подкольца. Каждый элемент $\alpha\in E$
задаёт гомоморфизм
\[
\varphi_\alpha\colon F[X]\to E,\qquad f(X)\mapsto f(\alpha).
\]
Возможны два случая.

\emph{Случай 1: $\ker\varphi_\alpha=0$}, т.\,е.\ для $f\in F[X]$
\[
f(\alpha)=0\ \Longrightarrow\ f=0\ \text{в}\ F[X].
\]
Тогда $\alpha$ называется \emph{трансцендентным} над $F$. Гомоморфизм $X\mapsto\alpha\colon F[X]\to F[\alpha]$~---
изоморфизм (сюръективен по \rf{1.21}, инъективен по условию), а если $E$~--- поле, то он продолжается до изоморфизма
полей частных $F(X)\to F(\alpha)$, $g/h\mapsto g(\alpha)/h(\alpha)$. Действительно, знаменатель $h(\alpha)\ne0$ при
$h\ne0$; отображение корректно (из $gh_1=g_1h$ следует $g(\alpha)h_1(\alpha)=g_1(\alpha)h(\alpha)$), является
гомоморфизмом полей, а значит, инъективно, и его образ~--- множество частных элементов $F[\alpha]$, т.\,е.\ $F(\alpha)$.
(Условие <<$E$~--- поле>> нужно лишь для того, чтобы частные имели смысл внутри $E$; для области целостности $E$ то же
верно в её поле частных.) В этом случае
$F[\alpha]$ не поле: $\alpha$ в нём не обратим, как не обратим $X$ в $F[X]$.

\emph{Случай 2: $\ker\varphi_\alpha\ne0$}, т.\,е.\ $g(\alpha)=0$ для некоторого ненулевого $g\in F[X]$. Тогда $\alpha$
называется \emph{алгебраическим} над $F$. Многочлены $g$ с $g(\alpha)=0$ образуют ненулевой идеал $\ker\varphi_\alpha$
кольца $F[X]$; по \rf{1.9} он порождается единственным унитарным многочленом $f$, а именно унитарным многочленом
наименьшей степени с $f(\alpha)=0$. Он называется \emph{минимальным многочленом} $\alpha$ над $F$. Он неприводим: если бы $f=gh$ с
$\deg g,\deg h<\deg f$, то $g(\alpha)h(\alpha)=0$, и по минимальности $g(\alpha)\ne0$, $h(\alpha)\ne0$~--- два ненулевых
элемента $E$ с нулевым произведением, что в области целостности невозможно.

\begin{proposition}[1.27c]
Пусть $\alpha\in E$ алгебраичен над $F$. Для унитарного $f\in F[X]$ следующие условия равносильны:
\begin{enumerate}[label=(\roman*),leftmargin=*]
\item $f(\alpha)=0$ и $f$ делит каждый $g\in F[X]$ с $g(\alpha)=0$;
\item $f$~--- унитарный многочлен наименьшей степени, для которого $f(\alpha)=0$;
\item $f$ неприводим и $f(\alpha)=0$.
\end{enumerate}
Каждое из них характеризует минимальный многочлен $\alpha$ над $F$. Кроме того, $g(X)\mapsto g(\alpha)$ индуцирует
изоморфизм $F[X]/(f)\to F[\alpha]$; в частности, $F[\alpha]$~--- поле, $(F[\alpha],\alpha)$~--- поле корня $f$, и
\[
[F[\alpha]:F]=\deg f.
\]
\end{proposition}

\begin{proof}
(i)$\Leftrightarrow$(ii): условие (i) означает, что $f$~--- унитарная образующая идеала $\ker\varphi_\alpha$, а по \rf{1.9}
это и есть унитарный многочлен наименьшей степени в этом идеале. Неприводимость такого $f$ доказана выше, т.\,е.\ (ii)$\Rightarrow$(iii).
(iii)$\Rightarrow$(i): если $f$ неприводим и $f(\alpha)=0$, то минимальный многочлен $f_0$ делит $f$, и, так как $f$
неприводим, а оба многочлена унитарны и $f_0\ne1$, имеем $f_0=f$.

Наконец, $\varphi_\alpha$ отображает $F[X]$ на $F[\alpha]$ с ядром $(f)$ и индуцирует изоморфизм $F[X]/(f)\to F[\alpha]$.
Слева поле степени $\deg f$ (\rf{1.25}), значит, и справа.
\end{proof}

Таким образом, для алгебраического $\alpha$ кольцо $F[\alpha]$ <<автоматически>> оказывается полем, и
$F(\alpha)=F[\alpha]$; это ключевое отличие от трансцендентного случая. Степень $[F[\alpha]:F]$ называется
\emph{степенью} алгебраического элемента $\alpha$ над $F$.

\begin{example}[1.28]
Пусть $\alpha\in\CC$ и $\alpha^3-3\alpha-1=0$. Многочлен $X^3-3X-1$ унитарен, неприводим и имеет корень $\alpha$, поэтому
по \rf{1.27c}(iii) он~--- минимальный многочлен $\alpha$ над $\Q$. Множество $\{1,\alpha,\alpha^2\}$~--- базис $\Q[\alpha]$
над $\Q$. Вычисления из \rf{1.27} (перенесённые в $\Q[\alpha]$ изоморфизмом $x\mapsto\alpha$ из \rf{1.27b}) показывают:
если $\beta=\alpha^4+2\alpha^3+3$, то $\beta=3\alpha^2+7\alpha+5$ и
\[
\beta^{-1}=\tfrac{7}{111}\alpha^2-\tfrac{26}{111}\alpha+\tfrac{28}{111}.
\]
\end{example}

\begin{remark}[1.29]
Системы компьютерной алгебры умеют вычислять в полях $\Q[a]$. Например, в PARI \texttt{factor(X\^{}4+4)} даёт разложение
\[
X^4+4=(X^2-2X+2)(X^2+2X+2)
\]
в $\Q[X]$ (ср.\ \rf{1.17c}). Команда \texttt{F=nfinit(a\^{}2+2*a+2)} задаёт числовое поле $\Q[a]$, порождённое корнем $a$ многочлена
$X^2+2X+2$, и \texttt{nffactor(F,x\^{}4+4)} возвращает разложение
\[
X^4+4=(X-a-2)(X-a)(X+a)(X+a+2)
\]
в $\Q[a][X]$. Проверим это вручную. Так как $a^2=-2a-2$, имеем
\[
(X-a)(X+a)=X^2-a^2=X^2+c,\qquad (X-a-2)(X+a+2)=X^2-(a+2)^2=X^2-c,\qquad c=2a+2,
\]
и произведение равно $X^4-c^2$, где $c^2=4(a^2+2a+1)=-4$. Итак, произведение равно $X^4+4$. (Над $\CC$ можно взять
$a=-1+i$; тогда множители~--- это $X\mp1\mp i$.)
\end{remark}

Расширение $E/F$ называется \emph{алгебраическим} (а $E$~--- алгебраическим над $F$), если каждый элемент $E$ алгебраичен
над $F$; в противном случае оно называется \emph{трансцендентным}. Таким образом, $E/F$ трансцендентно, если хотя бы один
элемент $E$ трансцендентен над $F$.

\begin{proposition}[1.30]
Пусть $E\supset F$~--- поля. Если $E/F$ конечно, то $E$ алгебраично и конечно порождено (как поле) над $F$. Обратно, если
$E$ порождено над $F$ конечным множеством алгебраических элементов, то $E/F$ конечно.
\end{proposition}

\begin{proof}
$\Rightarrow$. Сказать, что $\alpha\in E$ трансцендентен над $F$,~--- то же самое, что сказать, что степени
$1,\alpha,\alpha^2,\dots$ линейно независимы над $F$ (линейная зависимость $\sum a_i\alpha^i=0$~--- это в точности
ненулевой многочлен с корнем $\alpha$). Если $[E:F]=n<\infty$, то $n+1$ элементов $1,\alpha,\dots,\alpha^n$ линейно
зависимы, поэтому каждый $\alpha\in E$ алгебраичен над $F$.

Покажем, что $E$ конечно порождено над $F$. Если $E=F$, оно порождено пустым множеством. Иначе возьмём
$\alpha_1\in E\smallsetminus F$. Если $E\ne F[\alpha_1]$, возьмём $\alpha_2\in E\smallsetminus F[\alpha_1]$, и т.\,д. Все
кольца $F[\alpha_1,\dots,\alpha_k]$~--- подкольца конечного расширения, содержащие $F$, и потому поля (\rf{1.23}); в
частности, $F[\alpha_1,\dots,\alpha_k]=F(\alpha_1,\dots,\alpha_k)$. Процесс обрывается: каждое
$F[\alpha_1,\dots,\alpha_{k+1}]$ содержит $F[\alpha_1,\dots,\alpha_k]$ как \emph{собственное} $F$-подпространство
(элемент $\alpha_{k+1}$ в него не входит), а собственное подпространство конечномерного пространства имеет строго
меньшую размерность. Поэтому
\[
1<[F[\alpha_1]:F]<[F[\alpha_1,\alpha_2]:F]<\dots\le[E:F],
\]
и строго возрастающая последовательность натуральных чисел, ограниченная числом $[E:F]$, конечна; процесс
останавливается на некотором $n$, когда $E=F[\alpha_1,\dots,\alpha_n]$.

$\Leftarrow$. Пусть $E=F(\alpha_1,\dots,\alpha_n)$, где $\alpha_1,\dots,\alpha_n$ алгебраичны над $F$. Расширение
$F(\alpha_1)/F$ конечно: $F(\alpha_1)=F[\alpha_1]$ имеет степень $\deg f_1$, где $f_1$~--- минимальный многочлен $\alpha_1$
(\rf{1.27c}). Элемент $\alpha_2$ алгебраичен над $F$, а значит, и над $F(\alpha_1)$ (его минимальный многочлен над $F$~---
ненулевой многочлен с коэффициентами в $F\subset F(\alpha_1)$), поэтому
$F(\alpha_1,\alpha_2)=F(\alpha_1)(\alpha_2)$ конечно над $F(\alpha_1)$. По \rf{1.20} $F(\alpha_1,\alpha_2)$ конечно над $F$.
Продолжая по индукции (каждый этаж башни $F\subset F(\alpha_1)\subset F(\alpha_1,\alpha_2)\subset\dots\subset E$ конечен),
получаем, что $E/F$ конечно.
\end{proof}

Отсюда: \emph{если $\alpha_1,\dots,\alpha_n$ алгебраичны над $F$, то $F[\alpha_1,\dots,\alpha_n]=F(\alpha_1,\dots,\alpha_n)$, и
это поле конечно над $F$}. Действительно, $F(\alpha_1,\dots,\alpha_n)$ конечно над $F$ по \rf{1.30}, а
$F[\alpha_1,\dots,\alpha_n]$~--- его подкольцо, содержащее $F$, т.\,е.\ поле по \rf{1.23}.

Подчеркнём: в части $\Leftarrow$ существенна \emph{конечность} множества образующих. Алгебраическое расширение может
быть бесконечным (пример~--- поле всех алгебраических чисел, см.\ ниже \rf{1.47a}).

\begin{corollary}[1.31]
Пусть $L\supset E\supset F$~--- поля. Если $L$ алгебраично над $E$, а $E$ алгебраично над $F$, то $L$ алгебраично над $F$.
\end{corollary}

\begin{proof}
Пусть $\alpha\in L$. По условию $\alpha$~--- корень некоторого унитарного многочлена
\[
X^m+a_{m-1}X^{m-1}+\dots+a_0,\qquad a_i\in E.
\]
Элементы $a_0,\dots,a_{m-1}$ алгебраичны над $F$, поэтому по \rf{1.30} и замечанию после него $F[a_0,\dots,a_{m-1}]$~---
поле, конечное над $F$. Элемент $\alpha$ алгебраичен над $F[a_0,\dots,a_{m-1}]$ (написанный многочлен имеет коэффициенты в
этом поле), и по тому же замечанию $F[a_0,\dots,a_{m-1},\alpha]$~--- поле, конечное над $F[a_0,\dots,a_{m-1}]$. По \rf{1.20} поле $F[a_0,\dots,a_{m-1},\alpha]$ конечно над
$F$, и по \rf{1.30} его элемент $\alpha$ алгебраичен над $F$:
\[
F\ \underset{\text{конечно}}{\subset}\ F[a_0,\dots,a_{m-1}]\ \underset{\text{конечно}}{\subset}\ F[a_0,\dots,a_{m-1},\alpha]\ni\alpha.
\qedhere
\]
\end{proof}

\begin{proposition}[1.32]
Пусть $F$~--- поле, а $R$~--- область целостности, содержащая $F$ в качестве подкольца. Если $R$ порождается как
$F$-алгебра элементами, алгебраическими над $F$, то $R$~--- поле, алгебраическое над $F$.
\end{proposition}

Конечность $R$ над $F$ здесь не предполагается, и множество образующих может быть бесконечным.

\begin{proof}
Пусть сначала $R=F[\alpha_1,\dots,\alpha_n]$, где все $\alpha_i$ алгебраичны над $F$. Для каждого $i$ найдутся $m_i>0$ и
$c_{ij}\in F$ с
\[
\alpha_i^{m_i}=c_{i0}+c_{i1}\alpha_i+\dots+c_{i,m_i-1}\alpha_i^{m_i-1}
\]
(разделим на старший коэффициент многочлен с корнем $\alpha_i$). Это соотношение позволяет понижать степень $\alpha_i$ в
любом одночлене до значения $<m_i$, поэтому $R$ порождается как $F$-векторное пространство конечным набором элементов
\[
\alpha_1^{i_1}\cdots\alpha_n^{i_n},\qquad 0\le i_1<m_1,\ \dots,\ 0\le i_n<m_n.
\]
Значит, $R$~--- конечная $F$-алгебра без делителей нуля, т.\,е.\ поле (\rf{1.23}), конечное над $F$ и потому алгебраическое
(\rf{1.30}).

В общем случае каждый элемент $r\in R$ является многочленом от \emph{конечного} числа образующих (\rf{1.21}):
$r\in F[\alpha_1,\dots,\alpha_n]$ для некоторых алгебраических образующих $\alpha_1,\dots,\alpha_n$. По доказанному
$F[\alpha_1,\dots,\alpha_n]\subset R$~--- поле, алгебраическое над $F$; поэтому $r$ алгебраичен над $F$ и (при $r\ne0$)
обратим уже в $F[\alpha_1,\dots,\alpha_n]\subset R$.
\end{proof}

Например, кольцо $\Q[\sqrt2,\sqrt3,\sqrt5,\sqrt7,\dots]\subset\R$, порождённое квадратными корнями из всех простых чисел,
является полем, алгебраическим над $\Q$.

\subsection*{Примеры: квадратные корни и минимальные многочлены}

\begin{example}[1.32a: числа $1,\sqrt2,\sqrt3,\sqrt6$ линейно независимы над $\Q$]
Положим $K=\Q(\sqrt2,\sqrt3)\subset\R$. Докажем, что $[K:\Q]=4$ и что $1,\sqrt2,\sqrt3,\sqrt6$~--- базис $K$ над $\Q$.

\emph{Шаг 0: числа $\sqrt2$, $\sqrt3$, $\sqrt6$ иррациональны.} Многочлены $X^2-2$, $X^2-3$, $X^2-6$ неприводимы над $\Q$ по
критерию Эйзенштейна (\rf{1.16}) с $p=2$, $3$, $2$ соответственно (для $X^2-2$, $X^2-3$ это \rf{1.16a}(a)); неприводимый многочлен степени $2$ не имеет корней в
$\Q$. (Можно и напрямую: если $a^2=nb^2$ с целыми $a,b\ne0$, то каждое простое входит в $n$ в чётной степени, а у $n=2,3,6$
это не так.)

\emph{Шаг 1.} $X^2-2$~--- минимальный многочлен $\sqrt2$ над $\Q$ (\rf{1.27c}(iii)), поэтому $[\Q(\sqrt2):\Q]=2$ с базисом
$\{1,\sqrt2\}$.

\emph{Шаг 2: $\sqrt3\notin\Q(\sqrt2)$.} Пусть $\sqrt3=a+b\sqrt2$, $a,b\in\Q$. Возводя в квадрат, получаем
$3=a^2+2b^2+2ab\sqrt2$, и из линейной независимости $1,\sqrt2$ над $\Q$ следует $ab=0$. Если $b=0$, то $\sqrt3=a\in\Q$; если
$a=0$, то $\sqrt3=b\sqrt2$ и $\sqrt6=\sqrt2\cdot\sqrt3=2b\in\Q$. Оба случая противоречат шагу 0. Значит, $X^2-3$ не имеет
корней в $\Q(\sqrt2)$, т.\,е.\ неприводим над $\Q(\sqrt2)$ (см.\ замечание после \rf{1.12}), и $[K:\Q(\sqrt2)]=2$ с базисом $\{1,\sqrt3\}$.

\emph{Шаг 3.} По \rf{1.20} $[K:\Q]=2\cdot2=4$, и попарные произведения базисов $\{1,\sqrt2\}$ и $\{1,\sqrt3\}$, т.\,е.\
$1,\sqrt2,\sqrt3,\sqrt6$, образуют базис $K$ над $\Q$. В частности, они линейно независимы над $\Q$.

Тем же рассуждением доказывается, что $\sqrt2\notin\Q(\sqrt3)$ (если $\sqrt2=a+b\sqrt3$, то $ab=0$, и рационально $\sqrt2$
или $\sqrt6=3b$). Более общий факт (теорема Безиковича): квадратные корни из попарно различных бесквадратных натуральных чисел линейно
независимы над $\Q$; в частности, $[\Q(\sqrt{p_1},\dots,\sqrt{p_k}):\Q]=2^k$ для различных простых $p_1,\dots,p_k$.
Доказывается это с тем же приёмом <<возвести в квадрат и сравнить коэффициенты>>, но утверждение приходится усилить:
\emph{если бесквадратное $n>1$ имеет простой делитель вне $\{p_1,\dots,p_k\}$, то
$\sqrt n\notin\Q(\sqrt{p_1},\dots,\sqrt{p_k})$}. (Отсюда $\sqrt{p_k}\notin\Q(\sqrt{p_1},\dots,\sqrt{p_{k-1}})$ при каждом
$k$, и степень $2^k$ получается по \rf{1.20}, а с ней~--- базис из всех произведений
$\sqrt{p_{i_1}}\cdots\sqrt{p_{i_r}}$, $i_1<\dots<i_r$, откуда и линейная независимость корней из бесквадратных чисел.) Индукция по $k$; при $k=0$ утверждение означает иррациональность $\sqrt n$.
Пусть $K'=\Q(\sqrt{p_1},\dots,\sqrt{p_{k-1}})$, а бесквадратное $n$ имеет простой делитель $q\notin\{p_1,\dots,p_k\}$. По
предположению индукции $\sqrt{p_k}\notin K'$, т.\,е.\ $\{1,\sqrt{p_k}\}$~--- базис $K'(\sqrt{p_k})$ над $K'$. Если
$\sqrt n=a+b\sqrt{p_k}$, $a,b\in K'$, то $n=a^2+p_kb^2+2ab\sqrt{p_k}$, откуда $ab=0$. При $b=0$ получаем $\sqrt n\in K'$, а
при $a=0$ получаем $\sqrt{np_k}=bp_k\in K'$, т.\,е.\ $\sqrt m\in K'$ для бесквадратной части $m$ числа $np_k$
($np_k=mt^2$, $\sqrt m=\sqrt{np_k}/t$). Но и $n$, и $m$ делятся на $q$, и оба случая противоречат предположению индукции.

(С точки зрения теории Куммера степень $[\Q(\sqrt{a_1},\dots,\sqrt{a_k}):\Q]$ для $a_i\in\Q^\times$ равна порядку подгруппы,
порождённой классами $a_i$ в $\Q^\times/\Q^{\times2}$, см.\ \rf{5.31}(b) и \rf{5.31a}; например, $\Q(\sqrt2,\sqrt3,\sqrt6)=\Q(\sqrt2,\sqrt3)$
имеет степень $4$, а не $8$.)
\end{example}

\begin{example}[1.32b: минимальный многочлен $(2+\sqrt2)(3+\sqrt3)$]
Пусть $\gamma=(2+\sqrt2)(3+\sqrt3)=6+3\sqrt2+2\sqrt3+\sqrt6\in K=\Q(\sqrt2,\sqrt3)$. Докажем, что минимальный многочлен
$\gamma$ над $\Q$ равен
\[
P(X)=X^4-24X^3+144X^2-288X+144,
\]
и что $\Q(\gamma)=K$.

\emph{Как найти $P$.} Положим $\alpha=2+\sqrt2$, $\alpha'=2-\sqrt2$, $\beta=3+\sqrt3$, $\beta'=3-\sqrt3$; тогда
$\alpha+\alpha'=4$, $\alpha\alpha'=2$, $\beta+\beta'=6$, $\beta\beta'=6$, $\beta^2+\beta'^2=(\beta+\beta')^2-2\beta\beta'=24$.
Перемножим $X-\gamma$ с <<сопряжёнными>> множителями, сгруппировав их по два:
\[
(X-\alpha\beta)(X-\alpha'\beta)=X^2-(\alpha+\alpha')\beta X+\alpha\alpha'\beta^2=X^2-4\beta X+2\beta^2.
\]
Обратите внимание: $\beta$ входит в оба корня, поэтому в их произведении стоит $\beta^2$, а не $\beta$ (если написать
$2\beta$, получится неверный многочлен $X^4-24X^3+108X^2-96X+24$). Аналогично
$(X-\alpha\beta')(X-\alpha'\beta')=X^2-4\beta'X+2\beta'^2$, и
\begin{multline*}
(X^2-4\beta X+2\beta^2)(X^2-4\beta'X+2\beta'^2)\\
=X^4-4(\beta+\beta')X^3+\bigl(2(\beta^2+\beta'^2)+16\beta\beta'\bigr)X^2-8\beta\beta'(\beta+\beta')X+4(\beta\beta')^2,
\end{multline*}
что равно $X^4-24X^3+(48+96)X^2-288X+144=P(X)$. Итак, $P\in\Q[X]$, $P$ унитарен и $P(\gamma)=0$.

\emph{Почему $P$ минимален.} Минимальный многочлен $\gamma$ делит $P$, и достаточно доказать, что $[\Q(\gamma):\Q]=4$:
тогда минимальный многочлен имеет степень $4$ и совпадает с $P$ (оба унитарны). Так как $\Q(\gamma)\subset K$, по \rf{1.20}
степень $[\Q(\gamma):\Q]$ делит $[K:\Q]=4$, т.\,е.\ равна $1$, $2$ или $4$. Предположим, что она не больше $2$; тогда
$\gamma^2-s\gamma+p=0$ для некоторых $s,p\in\Q$. Вычислим в базисе $1,\sqrt2,\sqrt3,\sqrt6$ (\rf{1.32a}):
\[
\gamma^2=(\alpha\beta)^2=(6+4\sqrt2)(12+6\sqrt3)=72+48\sqrt2+36\sqrt3+24\sqrt6,
\]
\[
\gamma^2-s\gamma+p=(72-6s+p)+(48-3s)\sqrt2+(36-2s)\sqrt3+(24-s)\sqrt6.
\]
По единственности координат все четыре коэффициента равны нулю; но из $24-s=0$ следует $s=24$, и тогда $48-3s=-24\ne0$.
Противоречие. Значит, $[\Q(\gamma):\Q]=4=[K:\Q]$, откуда $\Q(\gamma)=K$ (подпространство полной размерности), и $P$~---
минимальный многочлен $\gamma$.

(Проверено на компьютере: \texttt{minimal\_polynomial((2+sqrt(2))*(3+sqrt(3)))} в sympy даёт $P$.)
\end{example}

\begin{example}[1.32c: минимальный многочлен $\sqrt{(2+\sqrt2)(3+\sqrt3)}$]
Пусть $\gamma$, $K$, $P$~--- как в \rf{1.32b}, и $\delta=\sqrt\gamma\in\R$. Тогда $\delta$~--- корень многочлена
\[
Q(X)=P(X^2)=X^8-24X^6+144X^4-288X^2+144.
\]
Докажем, что $Q$~--- минимальный многочлен $\delta$ над $\Q$, т.\,е.\ что $[\Q(\delta):\Q]=8$.

\emph{Сведение к вопросу о квадрате.} Так как $\gamma=\delta^2\in\Q(\delta)$ и $\Q(\gamma)=K$ (\rf{1.32b}), имеем
$K\subset\Q(\delta)$, т.\,е.\ $\Q(\delta)=K(\delta)$. По \rf{1.20}
\[
[\Q(\delta):\Q]=[K(\delta):K]\cdot[K:\Q]=4\,[K(\delta):K],
\]
а $[K(\delta):K]\in\{1,2\}$, так как $\delta$~--- корень $X^2-\gamma\in K[X]$. Итак,
\[
Q\ \text{минимален}\iff\delta\notin K\iff\gamma\ \text{не является квадратом в}\ K.
\]

\par\noindent\begin{minipage}{\linewidth}
\[
\begin{tikzpicture}[baseline]
\node (Q) at (0,0) {$\Q$};
\node (A) at (0,1.3) {$\Q(\sqrt3)$};
\node (K) at (0,2.6) {$K=\Q(\sqrt2,\sqrt3)=\Q(\gamma)$};
\node (L) at (0,3.9) {$\Q(\delta)=K(\delta)$};
\draw (Q)--(A)--(K)--(L);
\draw[decorate,decoration={brace,amplitude=4pt},cold] (2.5,2.45)--(2.5,0.15) node[midway,right=5pt]{$4$ (\rf{1.32a})};
\draw[decorate,decoration={brace,amplitude=4pt},cnew] (2.5,3.75)--(2.5,2.75) node[midway,right=5pt]{$2$, так как $\gamma\notin K^2$};
\draw[decorate,decoration={brace,amplitude=5pt},cnew] (-2.5,0.15)--(-2.5,3.75) node[midway,left=5pt]{$8$};
\end{tikzpicture}
\]
\figcap{Башня для $\delta=\sqrt\gamma$: степень нижней части известна, верхний этаж имеет степень $2$, как только $\gamma$~---
не квадрат в $K$.}
\end{minipage}\par\medskip

\emph{$\gamma$ не является квадратом в $K$.} Запишем элементы $K$ в виде $a+b\sqrt2$ с $a,b\in\Q(\sqrt3)$ (базис $\{1,\sqrt2\}$
над $\Q(\sqrt3)$: $[K:\Q(\sqrt3)]=4/2=2$, а $\sqrt2\notin\Q(\sqrt3)$ по \rf{1.32a}). Имеем
$\gamma=2(3+\sqrt3)+(3+\sqrt3)\sqrt2$. Пусть $\gamma=(a+b\sqrt2)^2=(a^2+2b^2)+2ab\sqrt2$. Сравнивая координаты,
\[
a^2+2b^2=2(3+\sqrt3),\qquad 2ab=3+\sqrt3.
\]
Тогда
\[
(a^2-2b^2)^2=(a^2+2b^2)^2-2(2ab)^2=4(3+\sqrt3)^2-2(3+\sqrt3)^2=2(3+\sqrt3)^2,
\]
т.\,е.\ $\bigl((a^2-2b^2)/(3+\sqrt3)\bigr)^2=2$, и $\sqrt2=\pm(a^2-2b^2)/(3+\sqrt3)\in\Q(\sqrt3)$. Это противоречит
\rf{1.32a}. Значит, $\gamma\notin K^2$, $[\Q(\delta):\Q]=8$, и $Q$~--- минимальный многочлен $\delta$ (он унитарен, имеет
корень $\delta$ и степень $8$). (Проверено в sympy: \texttt{minimal\_polynomial(sqrt((2+sqrt(2))*(3+sqrt(3))))} даёт $Q$.)
\end{example}

\begin{remark}[1.32d]
\begin{enumerate}[label=(\alph*),leftmargin=*]
\item Подстановка $X^2$ вместо $X$ в минимальный многочлен не всегда даёт минимальный многочлен квадратного корня: всё
решает, является ли элемент квадратом в порождённом им поле. Например, $3+2\sqrt2=(1+\sqrt2)^2$ имеет минимальный многочлен
$X^2-6X+1$, а у $\sqrt{3+2\sqrt2}=1+\sqrt2$ минимальный многочлен $X^2-2X-1$ степени $2$, а не $4$; соответственно
$X^4-6X^2+1=(X^2-2X-1)(X^2+2X-1)$ приводим.
\item Перемножение <<сопряжённых>> множителей, как в \rf{1.32b}, всегда даёт многочлен с корнем $\gamma$, но
минимальность нужно доказывать отдельно. Например, для $\sqrt2+\sqrt8$ произведение
$\prod(X\pm\sqrt2\pm\sqrt8)$ по всем знакам равно $X^4-20X^2+36=(X^2-2)(X^2-18)$, тогда как $\sqrt2+\sqrt8=3\sqrt2$ имеет
минимальный многочлен $X^2-18$. Здесь числа $1,\sqrt2,\sqrt8$ линейно зависимы ($\sqrt8=2\sqrt2$), и рассуждение
\rf{1.32b} не проходит.
\item Расширение $\Q(\delta)/\Q$ из \rf{1.32c}~--- известный пример: все восемь корней $Q$ лежат в $\Q(\delta)$ (например,
$\gamma\cdot(2-\sqrt2)(3+\sqrt3)=2(3+\sqrt3)^2$, так что $\sqrt{(2-\sqrt2)(3+\sqrt3)}=\pm\sqrt2\,(3+\sqrt3)/\delta\in\Q(\delta)$;
аналогично $\gamma\cdot(2+\sqrt2)(3-\sqrt3)=6(2+\sqrt2)^2$ и $\gamma\cdot(2-\sqrt2)(3-\sqrt3)=12$, так что
$\sqrt{(2+\sqrt2)(3-\sqrt3)}=\pm\sqrt6\,(2+\sqrt2)/\delta$ и $\sqrt{(2-\sqrt2)(3-\sqrt3)}=\pm2\sqrt3/\delta$ тоже лежат в
$\Q(\delta)$), и, как показывают методы теории Галуа (глава 3, см.\ упражнение 3-3), это расширение Галуа с группой
кватернионов $Q_8$.
\end{enumerate}
\end{remark}
